\chapter{Holomorphic-topological boundary conditions and 4d origins}
% label removed: ch:ht-boundary

\index{holomorphic-topological theory|textbf}

Every chiral algebra in Part~II arises as the boundary theory of a
holomorphic-topological field theory in the sense of Costello--Li.
The bar-cobar adjunction, from this perspective, encodes boundary-bulk
duality: the bar construction computes the open-string state space,
and bar-cobar inversion recovers the boundary algebra from bulk data.

\section{Mathematical relationships between frameworks}
% label removed: sec:precise-relationships-frameworks

\subsection{From 4D gauge theory to 2D chiral algebras}

\begin{theorem}[Costello--Li dimensional reduction; \ClaimStatusProvedElsewhere]% label removed: thm:costello-li-reduction
\index{Costello--Li construction}
\cite{CL16} Consider 4D $\mathcal{N}=2$ super Yang--Mills with gauge group $G$ 
on $\mathbb{C}^2$. After holomorphic-topological twist:
\begin{enumerate}
\item Fields become $\bar{\partial}$-closed differential forms with values in 
      $\mathfrak{g} \otimes \mathcal{O}_{\mathbb{C}^2}$
      
\item The action becomes BV-BRST exact up to a topological term: $S = \{Q_{\text{BRST}}, \Psi\} + S_{\mathrm{top}}$

\item Replacing $\mathbb{C}^2$ by $\mathbb{C} \times \mathbb{C}^*$ and taking $S^1$-equivariant reduction along $|w| = 1 \subset \mathbb{C}^*$
      produces a factorization algebra on $\mathbb{C}$
      
\item The resulting 2D theory has structure of a \emph{factorization algebra}
      $\mathcal{F}_G$, which is not a priori a chiral algebra
\end{enumerate}
\end{theorem}

\begin{remark}[Factorization algebras versus chiral algebras]% label removed: rem:fact-vs-chiral
The distinction is as follows (see \BDref{§2.3, §3.2}):

\emph{Factorization algebra} $\mathcal{F}$: assigns $\mathcal{F}(U)$ to every open set $U \subset X$, with multiplication maps $\mathcal{F}(U) \otimes \mathcal{F}(V) \to \mathcal{F}(U \sqcup V)$ for disjoint $U, V$ satisfying associativity (the factorization property over multiple disjoint opens).  Example: observables in any QFT.

\emph{Chiral algebra} $\mathcal{A}$: assigns a $\mathcal{D}_X$-module $\mathcal{A}_x$ at each point $x \in X$, with chiral product $j_*j^*(\mathcal{A} \boxtimes \mathcal{A}) \to \Delta_!\mathcal{A}$ having prescribed pole structure along the diagonal.  Example: vertex algebras, affine Kac--Moody algebras (Virasoro action is additional structure, not part of the definition).

\emph{Relationship} \cite{BD04,CG17}.
\[\text{Chiral algebras} \xleftrightarrow{\sim} \text{Factorization algebras on curves}\]
is an equivalence (Beilinson--Drinfeld \cite[Theorem~3.4.9]{BD04}).  In higher dimensions, chiral algebras embed into CG-style factorization algebras as those with $\mathcal{D}$-module structure and holomorphic dependence.
\end{remark}

\begin{conjecture}[When does CL produce chiral algebras?; \ClaimStatusConjectured]% label removed: conj:CL-produces-chiral
The Costello--Li construction produces a \emph{genuine chiral algebra} (not just
factorization algebra) if and only if:
\begin{enumerate}
\item The 4D theory has additional supersymmetry ensuring holomorphicity
\item The dimensional reduction preserves conformal symmetry
\item Central charge and anomaly terms satisfy consistency conditions
\end{enumerate}

For $\mathcal{N}=4$ SYM, this is established:
Theorem~\ref{thm:cl-n4-chirality} proves that
$\mathcal{F}_{G,\mathrm{Ad}}$ is genuinely chiral, producing the
affine Kac--Moody algebra $\widehat{\mathfrak{g}}_k$ at level~$k$
determined by the gauge coupling.  For $\mathcal{N}=2$ SYM with
general matter content, the precise conditions are conjectured in
Conjecture~\ref{conj:cl-general-chirality}.
\end{conjecture}

\begin{remark}[Evidence]
Chirality requires: (1) the twist preserves a holomorphic structure on the Coulomb branch \cite{Gai19}; (2) $T(z)$ survives the twist with vanishing anomalies; (3) the factorization algebra extends to a $\mathcal{D}$-module on $\mathrm{Ran}(X)$ (automatic for chiral algebras by \BDref{Chapter 3}, requires verification for twisted gauge theories).
\end{remark}

\begin{remark}[Scope]
Chirality is proved for $\mathcal{N}=4$ SYM (Theorem~\ref{thm:cl-n4-chirality}); the general case remains conjectural.  See \cite{CDG20,GKW24,Zeng23} for related work.  MC5 is proved at all genera in Volume~I; the remaining obstruction is the chirality verification for general twisted gauge theories.
\end{remark}

\subsection{Costello--Li conditions: from factorization to chiral}

\begin{theorem}[CL chirality for \texorpdfstring{$\mathcal{N}=4$}{N=4}; \ClaimStatusProvedElsewhere]
% label removed: thm:cl-n4-chirality
\index{Costello--Li construction!chirality criterion}
\cite{CL16}
Let $G$ be a reductive group and consider 4d $\mathcal{N}=2$ gauge theory
with gauge group $G$ and matter in a representation $R$ on $\Sigma \times \mathbb{C}$.
After holomorphic twist, the theory produces a factorization algebra
$\mathcal{F}_{G,R}$ on $\Sigma$.  When the theory has $\mathcal{N}=4$
supersymmetry (equivalently, $R = \mathrm{Ad}(G)$), the factorization algebra
is strictly chiral: $\mathcal{F}_{G,\mathrm{Ad}}$ satisfies the chiral algebra
axioms of Beilinson--Drinfeld.  The resulting chiral algebra is the affine
Kac--Moody algebra $\widehat{\fg}_k$ at level $k$ determined by the gauge coupling.
\end{theorem}

\begin{conjecture}[CL chirality conditions for general \texorpdfstring{$R$}{R}; \ClaimStatusConjectured]
% label removed: conj:cl-general-chirality
\index{Costello--Li construction!$\mathcal{N}=2$ conditions}
For $\mathcal{N}=2$ gauge theory with gauge group $G$ and matter content $R$,
the factorization algebra $\mathcal{F}_{G,R}$ is chiral provided:
\begin{enumerate}[label=\textup{(\roman*)}]
\item $R$ is self-dual: $R \cong R^*$ as a $G$-representation.
\item The one-loop anomaly $c_1(R) \in H^2(\Sigma, \mathbb{Z})$ vanishes.
\item The matter content is balanced: $\chi(\Sigma, R) = 0$ for all genera.
\end{enumerate}
When these conditions hold, the resulting chiral algebra $\mathcal{A}_{\mathrm{CL}}(G, R)$
on $\Sigma$ has:
\begin{itemize}
\item Central charge $c = \dim(G) \cdot (1 - h^\vee/k) + \tfrac{1}{2}\dim(R)$,
  where $k$ is the level and $h^\vee$ is the dual Coxeter number.
\item Koszul dual $\mathcal{A}_{\mathrm{CL}}(G, R)^! = \mathcal{A}_{\mathrm{CL}}(G^\vee, R^\vee)$,
  where $G^\vee$ is the Langlands dual group and $R^\vee$ the dual representation.
\item Bar complex $\barB(\mathcal{A}_{\mathrm{CL}})$ computing the derived moduli
  of solutions to the BPS equations on $\Sigma \times \mathbb{C}$.
\end{itemize}
For $G = \mathrm{GL}(N)$ and $R = \mathrm{Ad}$, the chiral algebra is the
$\mathcal{W}_N$-algebra, and the instanton partition function on $\mathbb{C}^2$
equals the genus-$0$ bar complex evaluation; this is the AGT correspondence
(Theorem~\ref{thm:agt}).
\end{conjecture}

\begin{remark}[Scope]
Conditions~(i)--(iii) are established for $\mathcal{N}=4$; the general case is conjectural.
\end{remark}

\subsection{Paquette--Williams boundary vertex algebras}

\begin{theorem}[Paquette--Williams 2022; \ClaimStatusProvedElsewhere]% label removed: thm:paquette-williams-boundaries
\cite{PW22} Consider holomorphic-topological 4D $\mathcal{N}=2$ gauge theory on 
$\mathbb{C}^2$ with boundary at $z_2 = 0$. Then:
\begin{enumerate}
\item Boundary conditions $\mathcal{B}$ correspond to Lagrangian submanifolds in 
      Coulomb branch $\mathcal{M}_C$
      
\item The boundary supports a \emph{vertex operator algebra} (VOA) $V_{\mathcal{B}}$

\item This VOA is the quantization of the symplectic reduction of $\mathcal{M}_C$ 
      at the boundary
      
\item The VOA $V_{\mathcal{B}}$ has a \emph{chiral envelope} $\mathcal{A}_{\mathcal{B}}$,
      which is a chiral algebra in the BD sense
\end{enumerate}
\end{theorem}

\begin{remark}[Bar-cobar duality for boundary VOAs]% label removed: rem:PW-connection
Paquette--Williams produce vertex algebras, which are algebraic objects (Frenkel--Ben-Zvi
\cite{FBZ04}). The bar-cobar framework of Part~I applies to their \emph{chiral envelopes}, which are
geometric objects ($\mathcal{D}$-modules).

\begin{center}
\begin{tikzcd}
\text{Boundary VOA } V_{\mathcal{B}} \ar[r, "\text{chiral envelope}"] \ar[d, "\text{PW}"]
& \text{Chiral algebra } \mathcal{A}_{\mathcal{B}} \ar[d, "\text{bar-cobar}"] \\
\text{4D HT gauge theory} \ar[r, "\text{CL reduction}"]
& \text{2D factorization algebra}
\end{tikzcd}
\end{center}

The vertex algebra $V_{\mathcal{B}}$ contains the algebraic information (OPE, modes, 
etc.), while the chiral algebra $\mathcal{A}_{\mathcal{B}}$ contains the geometric 
information (configuration spaces, $\mathcal{D}$-modules, sheaf cohomology).

Bar-cobar duality (Part~I) applies to $\mathcal{A}_{\mathcal{B}}$: the bar construction $\bar{B}(\mathcal{A}_{\mathcal{B}})$ \emph{is} the factorization coalgebra on $\Ran(X)$ that encodes PW boundary data (Convention~\ref*{V1-conv:bar-coalgebra-identity}).  Verdier duality $\mathbb{D}_{\Ran}$ converts this coalgebra into a factorization algebra whose cohomology on the Koszul locus is $\mathcal{A}_{\mathcal{B}}^!$, identifying the dual boundary condition.  Separately, bar-cobar inversion $\Omega\bar{B}(\mathcal{A}_{\mathcal{B}}) \xrightarrow{\sim} \mathcal{A}_{\mathcal{B}}$ recovers the original algebra (a different operation).
\end{remark}

\section{From 4d SYM to holomorphic Chern--Simons}

\subsection{The holomorphic twist and localization}

\begin{definition}[Holomorphic twist of 4d \texorpdfstring{$\mathcal{N}=4$}{N=4} SYM]
Start with $\mathcal{N}=4$ super Yang--Mills in 4d with gauge group $G$, viewed as an $\mathcal{N}=2$ theory by choosing a preferred $\mathcal{N}=2$ subalgebra (Costello~\cite{costello-renormalization}).  The field content before twisting:
\begin{itemize}
\item Vector multiplet: $A_\mu$ (gauge field), $\Phi^I$ (6 scalars), $\lambda, 
\bar{\lambda}$ (fermions)
\item Supersymmetry: 16 supercharges transforming under $\text{Spin}(6)_R$
\end{itemize}

The \emph{holomorphic twist} (not to be confused with the fully topological Kapustin--Witten A-twist of $\mathcal{N}=4$):
\begin{enumerate}
\item Decompose $\mathbb{R}^4 = \mathbb{C} \times \mathbb{C}$ with coordinates 
$(z, w)$
\item Twist the Lorentz group: $\text{SO}(4) \to \text{SO}(2)_{\text{hol}} \times 
\text{SO}(2)_{\text{top}}$
\item Mix with R-symmetry to make a supercharge $Q$ scalar
\item Result: Theory is holomorphic in $z$, topological in $w$
\end{enumerate}
\end{definition}

\begin{theorem}[Localization to holomorphic data; \ClaimStatusProvedElsewhere]% label removed: thm:ht-localization
After the holomorphic twist \cite{costello-renormalization}, the path integral localizes to:
\[Z = \int_{[\bar{\partial}_A = 0]} \mathcal{O}(\text{fields}) \cdot e^{-S_{\text{inst}}}\]

where $\bar{\partial}_A = 0$ means:
\begin{itemize}
\item $A$ is a holomorphic connection
\item Scalars $\Phi$ satisfy holomorphic moment map equation
\end{itemize}

This is the moduli space of holomorphic $G$-bundles with Higgs field.
\end{theorem}

\begin{proof}[Sketch]
The twisted action decomposes as $S_{\mathrm{twisted}} = \{Q, V\} + S_0$ where $Q$ is the scalar supercharge, $V$ is the gauge fermion, and $S_0$ is metric-independent.  Since $\{Q, V\} \geq 0$, the path integral localizes to
\[\mathcal{M}_Q = \{F_{A}^{(0,2)} = 0,\; F_A^{(1,1)} + [\Phi, \Phi^*] = 0,\; \bar{\partial}_A \Phi = 0\},\]
i.e., Hitchin's self-duality equations in the holomorphic gauge.
\end{proof}

\subsection{Holomorphic Chern--Simons as effective theory}

\begin{definition}[Holomorphic Chern--Simons action]
\index{holomorphic Chern--Simons}
On a Calabi--Yau 3-fold $X$ with holomorphic volume form $\Omega \in \Omega^{3,0}(X)$, the
holomorphic Chern--Simons action is:
\[S_{\text{HCS}}[A] = \frac{1}{2}\int_{X} \Omega \wedge \operatorname{Tr}\left(A
\wedge \bar{\partial} A + \frac{2}{3} A \wedge A \wedge A\right)\]

where $A \in \Omega^{0,1}(X, \mathfrak{g}_{\mathbb{C}})$.  Since $A \in \Omega^{0,1}$ and $\bar\partial A \in \Omega^{0,2}$, we have $A \wedge \bar\partial A \in \Omega^{0,3}$ and $A^3 \in \Omega^{0,3}$, so $\Omega \wedge \text{Tr}(\cdots)$ is a $(3,3)$-form, hence integrable over~$X$.  The factor $\frac{1}{2}$ follows the convention of Costello~\cite{costello-renormalization}.
\end{definition}

\begin{theorem}[HCS from dimensional reduction; \ClaimStatusProvedElsewhere]% label removed: thm:hcs-dimensional-reduction
Holomorphic Chern--Simons \cite{costello-renormalization} arises from the holomorphic twist of 4d SYM by:
\begin{enumerate}
\item The holomorphic twist on $\mathbb{C}^2$ yields a theory holomorphic in $z$, topological in~$w$
\item Compactify the topological direction $w$ on $S^1$ to produce a 3d theory on $\mathbb{C} \times \mathbb{R}$
\item This 3d theory is holomorphic Chern--Simons
\end{enumerate}

The 4d holomorphic twist yields a holomorphic-topological theory on $\mathbb{C}^2$. Compactifying one direction produces a 3d theory; the holomorphic volume form on the resulting CY 3-fold $X$ is a $(3,0)$-form $\Omega \in \Omega^{3,0}(X)$.  Locally, taking $X = \mathbb{C}^3$ with coordinates $(z, w, u)$:
\[\Omega = dz \wedge dw \wedge du\]
\end{theorem}

\section{Boundary conditions and chiral operads}

\subsection{The deformed conifold geometry}

\begin{example}[Costello--Gaiotto holography setup]
The canonical example is the deformed conifold:
\[X = \{u_1 w_2 - u_2 w_1 = N\} \subset \mathbb{C}^4\]

This space has:
\begin{itemize}
\item Calabi--Yau 3-form: $\Omega|_X = \mathrm{Res}_{f=N}\frac{du_1 \wedge du_2 \wedge dw_1
\wedge dw_2}{f - N}$ (Poincar\'e residue of the ambient 4-form along $X = \{f = N\}$, yielding a holomorphic 3-form on $X$)
\item $SL_2(\mathbb{C})$ isometry: Acting on $(u, w)$ coordinates
\item Asymptotic boundary: $\mathbb{CP}^1 \times \mathbb{CP}^1$ as $|u|, |w| 
\to \infty$
\item Pole structure: $\Omega$ has cubic pole along the boundary divisor
\end{itemize}
\end{example}


\subsection{HT boundary conditions}

\begin{definition}[Holomorphic-topological boundary condition]
For HCS on $X$ with boundary compactification $\bar{X}$ and boundary divisor 
$D = \bar{X} \setminus X$:

A \emph{holomorphic-topological boundary condition} specifies:
\begin{enumerate}
\item Extension: Fields extend to $\bar{X}$ as holomorphic sections
\item Vanishing order: $A \in H^0(\bar{X}, \Omega^{0,1}(\mathcal{O}(-D)))$ (simple zero at $D$)
\item Behavior: As approaching $D$, $A \sim (\text{distance}) \cdot (\text{smooth})$
\end{enumerate}
\end{definition}

\begin{theorem}[Boundary chiral algebra; \ClaimStatusProvedElsewhere]
% label removed: thm:ht-boundary-chiral-algebra
\cite{CDG20}
An HT boundary condition supports a chiral algebra $\mathcal{A}_{\text{bdy}}$ whose:
\begin{enumerate}
\item \emph{Generators.} Boundary local operators $\mathcal{O}(z)$ for $z \in D$
\item \emph{OPE.} Determined by bulk path integral with boundary insertions
\[\mathcal{O}_1(z) \mathcal{O}_2(w) = \sum_k C_{12}^k(z-w) \mathcal{O}_k(w)\]
\item \emph{Factorization.} Extends to factorization algebra on $D$
\end{enumerate}
\end{theorem}

\begin{proof}[Sketch]
Boundary operators $\mathcal{O}(z) = \lim_{\epsilon \to 0} \mathrm{Tr}(A(z + \epsilon n) \cdots)$ are well-defined by the simple zero condition.  The OPE arises from short-distance singularities of the HCS path integral with boundary insertions.  Locality (annular convergence), associativity (path integral composition), and skyscraper support on~$D$ give the BD chiral algebra axioms.
\end{proof}

\subsection{Chiral operad action}

\begin{conjecture}[Chiral operad from HCS; \ClaimStatusConjectured]
% label removed: conj:chiral-operad-hcs
The holomorphic Chern--Simons theory defines a chiral operad $\mathcal{P}_{\text{HCS}}$ 
acting on boundary chiral algebras.

The operad operations:
\[\mathcal{P}_{\text{HCS}}(n) = \text{Obs}(X; D \times \overline{C}_n(D))\]

are observables on $X$ with marked points on the boundary.
\end{conjecture}

\begin{remark}[Scope]
The operad structure is predicted by \cite{CG17}; the explicit identification with HCS observables requires perturbative quantization.
\end{remark}

\begin{example}[Kac--Moody from gauge HCS]
For $G = SU(N)$ holomorphic Chern--Simons:
\[\mathcal{A}_{\text{bdy}} = \widehat{\mathfrak{sl}}_N\]

the affine Kac--Moody algebra at level $k$ (determined by HCS coupling).

The boundary currents:
\[J^a(z) = \lim_{\epsilon \to 0} \text{Tr}(T^a A(z + \epsilon n))\]

satisfy the Kac--Moody OPE:
\[J^a(z) J^b(w) \sim \frac{k \delta^{ab}}{(z-w)^2} +
\frac{f^{abc} J^c(w)}{z-w}\]

This is derivable from the HCS path integral.
\end{example}

\section{Open-closed correspondence as bar-cobar duality}

\subsection{Open and closed sectors}

\begin{theorem}[Open-string bar identification; \ClaimStatusProvedHere]
% label removed: thm:open-string-bar
\index{open-closed correspondence!open sector}
Let $\mathcal{A}_{\mathrm{bdy}}$ be the boundary chiral algebra of
a holomorphic-topological theory with boundary condition supporting
a Kac--Moody or W-algebra vertex algebra.  Then the bar complex
$\barB^{\mathrm{ch}}(\mathcal{A}_{\mathrm{bdy}})$ computes the
open-string state space: there is a quasi-isomorphism
\[
\barB^{\mathrm{ch}}(\mathcal{A}_{\mathrm{bdy}})
\simeq
C^{\frac{\infty}{2}+\bullet}(\mathcal{A}_{\mathrm{bdy}})
\]
to the semi-infinite (BRST) complex of~$\mathcal{A}_{\mathrm{bdy}}$.
Moreover, Verdier duality on $\operatorname{Ran}(X)$ converts the
bar coalgebra $\barB^{\mathrm{ch}}(\cA_{\mathrm{bdy}})$ into a
factorization algebra whose cohomology, on the Koszul locus, is the
Koszul dual~$\cA_{\mathrm{bdy}}^!$
\textup{(}Theorem~\textup{\ref*{V1-thm:verdier-bar-cobar})}.
This exchanges boundary and bulk data: the coalgebra encodes the
boundary; its Verdier dual algebra encodes the bulk.
\end{theorem}

\begin{proof}
For Kac--Moody boundary conditions
$\mathcal{A}_{\mathrm{bdy}} = \widehat{\fg}_k$ ($k \neq -h^\vee$),
Theorem~\ref*{V1-thm:bar-semi-infinite-km} provides the quasi-isomorphism
$\barB^{\mathrm{ch}}(\widehat{\fg}_k) \simeq C^{\frac{\infty}{2}+\bullet}(\widehat{\fg}_k)$
via the PBW filtration and Arkhipov's functor.
For W-algebra boundary conditions
$\mathcal{A}_{\mathrm{bdy}} = \mathcal{W}_k(\fg)$,
Theorem~\ref*{V1-thm:bar-semi-infinite-w} gives the analogous
identification through Drinfeld--Sokolov reduction.
In both cases, the semi-infinite complex is the mathematical
formulation of the open-string BRST complex: its cohomology
computes the physical state space of open strings ending on
the holomorphic-topological boundary condition.

For the Verdier duality statement:
Theorem~\ref*{V1-thm:verdier-bar-cobar} shows that $\mathbb{D}_{\operatorname{Ran}}$
maps the bar coalgebra to a factorization algebra.  On the Koszul
locus (which includes Kac--Moody and W-algebra boundary conditions
at generic level), the cohomology of the Verdier dual is $\cA_{\mathrm{bdy}}^!$.
Note: Verdier duality converts the coalgebra $\barB(\cA)$ to an
algebra; this is \emph{not} the same operation as the cobar
functor $\Omega$, which is the left adjoint to~$\barB$.
\end{proof}

\begin{conjecture}[Closed-string cobar identification; \ClaimStatusConjectured]
% label removed: conj:closed-string-cobar
\index{open-closed correspondence!closed sector}
Let $\mathcal{C}_{\mathrm{bulk}}$ be the bulk chiral coalgebra of
a holomorphic-topological string theory.  Then:
\begin{enumerate}
\item \emph{Closed strings.} The cobar complex
  $\Omega^{\mathrm{ch}}(\mathcal{C}_{\mathrm{bulk}})$ computes the
  closed-string field theory state space in the sense of
  Zwiebach~\cite{Zwi93}.
\item \emph{Open-closed duality.} The bar-cobar adjunction
  should furnish the full open-closed correspondence: the natural map
  $\barB(\mathcal{A}_{\mathrm{bdy}}) \to
  \Omega(\mathcal{C}_{\mathrm{bulk}})$ is conjecturally the open-closed string
  field theory map.
\end{enumerate}
\end{conjecture}

\begin{remark}[Evidence]
Open sector: boundary vertex operators are elements of $\mathcal{A}_{\mathrm{bdy}}$, and off-shell amplitudes lie in $\barB^{\mathrm{ch}}(\mathcal{A}_{\mathrm{bdy}})$ (proved).  Closed sector: bulk vertex operators give distribution-valued correlation functions on open configuration spaces (conjectural).  The open-closed map $\barB(\mathcal{A}_{\mathrm{bdy}}) \to \Omega(\mathcal{C}_{\mathrm{bulk}})$ sends boundary residues to bulk distributions via Stokes' theorem.
\end{remark}

\begin{remark}[Scope]
Open-string side proved (Theorems~\ref*{V1-thm:bar-semi-infinite-km}, \ref*{V1-thm:bar-semi-infinite-w}).  MC4 and MC5 are both proved in Volume~I; the closed-string identification is now accessible via the all-genera bar-cobar framework.
\end{remark}

\subsection{Factorization and dimensional reduction}

\begin{conjecture}[Factorization along dimension tower; \ClaimStatusConjectured]% label removed: conj:factorization-dim-tower
The dimensional reduction sequence:
\[\text{4d SYM} \xrightarrow{\text{HT twist + reduce}}
\text{3d HCS} \xrightarrow{\text{boundary}}
\text{2d chiral algebra} \xrightarrow{\text{defect}}
\text{1d quantum mechanics}\]

is governed by iterated bar-cobar constructions at each level.
\end{conjecture}

\begin{remark}[Scope]
Individual levels are established \cite{BD04,CG17,CWY18}; MC5 is proved at all genera in Volume~I, so the bar-cobar side of the iterated identification is now unconditional.
\end{remark}

\begin{example}[Explicit tower for \texorpdfstring{$\mathcal{N}=4$}{N=4} SYM]
$\mathcal{N}=4$ SYM with gauge group $G$ on $\mathbb{R}^4$
$\xrightarrow{\text{HT twist + reduce}}$
HCS on $\mathbb{C} \times \mathbb{R}_{\geq 0}$
$\xrightarrow{\text{boundary}}$
$\widehat{\mathfrak{g}}_k$ on $\mathbb{C}$
$\xrightarrow{\text{defect}}$
quantum integrable system (Toda for $G = SU(N)$).
\end{example}

\section{$\mathcal{W}$-algebras from Hitchin moduli}

\subsection{The Higgs branch and Hitchin system}

\begin{definition}[Hitchin moduli space]
For a Riemann surface $\Sigma_g$ of genus $g$ and gauge group $G$, the Hitchin 
moduli space is:
\[\mathcal{M}_{\text{Hit}}(\Sigma_g, G) = \{(E, \Phi) : \bar{\partial}_E \Phi = 0\} / 
\sim\]

where $E \to \Sigma_g$ is a holomorphic $G$-bundle, $\Phi \in H^0(\Sigma_g, \operatorname{End}(E) \otimes K_{\Sigma_g})$ is the Higgs field, and $\sim$ denotes gauge equivalence.
\end{definition}

\begin{proposition}[Virasoro from \texorpdfstring{$SL_2$}{SL_2}-Hitchin; \ClaimStatusProvedElsewhere]% label removed: prop:virasoro-from-hitchin
For $G = SL_2$, the chiral algebra of local operators on the Hitchin moduli space $\mathcal{M}_{\mathrm{Hit}}(\Sigma_g, SL_2)$ is the Virasoro algebra:
\[\mathcal{A}_{\mathrm{local}}(\mathcal{M}_{\mathrm{Hit}}(\Sigma_g, SL_2)) \cong \mathrm{Vir}_c.\]
\end{proposition}

\begin{proof}
The Hitchin system for $SL_2$ is the cotangent bundle $T^*\mathrm{Bun}_{SL_2}(\Sigma_g)$; quantization of the coordinate ring via the Beilinson--Drinfeld construction \cite{BD04} produces the Virasoro vertex algebra at the level determined by the symplectic structure.
\end{proof}

\begin{conjecture}[W-algebra from Hitchin; \ClaimStatusConjectured]
% label removed: conj:w-algebra-hitchin
For general $G$, the chiral algebra of local operators on $\mathcal{M}_{\mathrm{Hit}}(\Sigma_g, G)$ is
\[
\mathcal{A}_{\mathrm{local}}(\mathcal{M}_{\mathrm{Hit}}) \cong \mathcal{W}(G),
\]
the W-algebra associated to~$G$.  For $G = SL_N$, this is the
$\mathcal{W}_N$ algebra with generators
$T(z)$, $W^{(3)}(z)$, \ldots, $W^{(N)}(z)$
of conformal weights $2, 3, \ldots, N$.
\end{conjecture}

\begin{remark}[Scope]
The $SL_2$ case is Beilinson--Drinfeld \cite{BD04}; the general case requires AGT localization.  Partial results: Schiffmann--Vasserot \cite{SV13}, Maulik--Okounkov \cite{MO19}, Arakawa \cite{Ara07}.
\end{remark}

\begin{proof}[Via AGT correspondence]
Place 4d $\mathcal{N}=2$ gauge theory with gauge group $G$ on $\mathbb{R}^2_\epsilon \times \Sigma_g$ with $\Omega$-background parameters $(\epsilon_1, \epsilon_2)$.  The $\Omega$-background localizes the path integral to $\mathcal{M}_{\mathrm{Hit}}$.  In the Nekrasov limit $\epsilon_2 \to 0$:
\[Z_{4d}|_{\epsilon_2 \to 0} = Z_{\mathcal{W}}[\Sigma_g],\]
the $\mathcal{W}(G)$ partition function on $\Sigma_g$.  The operator correspondence (4d line operators $\leftrightarrow$ 2d vertex operators \cite{DNP25}, 4d surface operators $\leftrightarrow$ 2d extended operators) identifies local operators on $\mathcal{M}_{\mathrm{Hit}}$ with generators of $\mathcal{W}(G)$.
\end{proof}

\subsection{Bar-cobar for $\mathcal{W}$-algebras}

\begin{theorem}[\texorpdfstring{$\mathcal{W}$}{W}-algebra bar complex; \ClaimStatusProvedHere]
% label removed: thm:w-algebra-bar-complex
For $\mathcal{W}_N$, the geometric bar complex:
\[\bar{B}^{\text{ch}}(\mathcal{W}_N) =
\bigoplus_{n \geq 0} \Gamma\!\left(\overline{C}_{n+1}(X),\,
  j_*j^*\mathcal{W}_N^{\boxtimes(n+1)} \otimes \Omega^n_{\log D}\right)\]

has three proved outputs:
\begin{enumerate}
\item \textup{[\ClaimStatusProvedHere]}
\emph{Conformal blocks.}
The bar-complex cohomology
$H^0(\bar{B}^{\mathrm{ch}}(\mathcal{W}_N))$ computes conformal blocks.
This is the Part~I bar-complex/chiral-homology correspondence applied
to the specific chiral algebra $\mathcal{W}_N$.

\item \textup{[\ClaimStatusProvedHere]}
\emph{Fusion rules.} For rational $\mathcal{W}_N$
\textup{(}$C_2$-cofinite, Arakawa \cite{Ara07}\textup{)},
the bar differential with module insertions at marked
points encodes fusion: $\mathrm{Ext}^1_{\mathcal{W}_N}(M_i
\otimes M_j, M_k)$ is computed by the bar complex
$\barB(\mathcal{W}_N; M_i, M_j, M_k)$, and fusion
multiplicities $N_{ij}^k = \dim
\mathrm{Hom}(M_k, M_i \boxtimes M_j)$ are extracted via
residues at collision divisors.

\item \textup{[\ClaimStatusProvedHere]}
\emph{Modular functors.} The genus-$g$ bar complex
$\barB^{(g)}(\mathcal{W}_N)$ defines a vector bundle
over $\overline{\mathcal{M}}_{g,n}$ whose fibers are
conformal block spaces; the modular operad structure
\textup{(}Theorem~\textup{\ref*{V1-thm:genus-induction-strict})}
governs the genus dependence, and
$\barB^{(g)}$ computes chiral homology
\textup{(}Theorem~\textup{\ref*{V1-thm:genus-g-chiral-homology})}.
\end{enumerate}
\end{theorem}

\begin{remark}[Scope]
All three items are proved.  Item~(2) uses the Huang--Lepowsky--Zhang tensor product theory \cite{HLZ} for $C_2$-cofinite algebras.  Explicit Virasoro modular functor computations are in Theorem~\ref*{V1-thm:chain-modular-functor}.
\end{remark}

\begin{example}[Virasoro = \texorpdfstring{$\mathcal{W}_2$}{W_2}]
For $G = SL_2$, $\mathcal{W}_2 = \text{Vir}$ is the Virasoro algebra.

The bar complex element:
\[\omega = T(z_1) \otimes T(z_2) \otimes \cdots \otimes T(z_n) \otimes 
\bigwedge_{i<j} \eta_{ij}^{k_{ij}}\]

represents an off-shell correlator of stress tensors.

The bar differential:
\begin{itemize}
\item $d_{\text{res}}$: Extracts OPE $T(z_1)T(z_2) = \frac{c/2}{(z_1-z_2)^4} + \cdots$
\item $d_{\text{strat}}$: Accounts for degeneration of moduli space
\item $d_{\text{int}}$: Implements Ward identities
\end{itemize}

On-shell correlators (physical observables) are in:
\[H^0(\bar{B}^{\text{ch}}(\text{Vir}))\]
\end{example}

\section{Quantization and loop corrections}

\subsection{Classical vs.\ quantum chiral algebras}

\begin{definition}[Quantum correction parameter]
In the reduction from 4d to 2d, there is a natural parameter:
\[\hbar = \epsilon_1\]

This is the $\Omega$-background parameter, which becomes Planck's constant in 
the reduced theory.

Classical limit: $\hbar \to 0$ (or $\epsilon_1 \to 0$)

Quantum theory: $\hbar$ finite
\end{definition}

\begin{theorem}[Genus-graded bar complex; \ClaimStatusProvedHere]
% label removed: thm:genus-graded-bar
The full bar complex has a genus-graded deformation
\[\bar{B}^{\mathrm{ch}}_\hbar(\mathcal{A}) = \bar{B}^{\mathrm{ch}}(\mathcal{A})[[\hbar]]\]
with differential
\[d_\hbar = d_0 + \hbar d_1 + \hbar^2 d_2 + \cdots\]
where $d_g$ (the genus-$g$ component of $\Dg{\bullet}$, Convention~\ref*{conv:higher-genus-differentials}) is the genus-$g$ contribution from the modular operad decomposition.
The nilpotency $(d_\hbar)^2 = 0$ holds as a formal power series identity.
\end{theorem}

\begin{proof}
By Theorem~\ref*{V1-thm:genus-induction-strict}, the full bar complex
$\bar{B}^{\mathrm{full}}(\mathcal{A}) = \bigoplus_g \bar{B}^{(g)}(\mathcal{A})$
carries a differential whose genus-$g$ component $d_g$ (in the sense of Convention~\ref*{conv:higher-genus-differentials}) is assembled from
the modular operad structure.  Setting $\hbar^g$ as the genus bookkeeping
parameter, the relation $(d_\hbar)^2 = 0$ is the genus-by-genus
identity $\sum_{g_1 + g_2 = g} d_{g_1} \circ d_{g_2} = 0$
for each $g$, which is proved in Theorem~\ref*{V1-thm:genus-induction-strict}.
\end{proof}

\begin{remark}[Physics interpretation]
The nilpotency $(d_\hbar)^2 = 0$ is the Maurer--Cartan equation $\ell_1(\Theta) + \tfrac{1}{2}\ell_2(\Theta, \Theta) = 0$ of Theorem~\ref*{V1-thm:explicit-theta}, expanded genus-by-genus.  In the Costello--Li framework, $\hbar = \epsilon_1$ and $d_k$ is the $k$-loop correction; this identification requires gauge-theoretic input.
\end{remark}

\begin{example}[Two classical limits of Virasoro]
The Lie-algebraic limit $c = 0$ gives the Witt algebra (no central extension, $\Theta_\cA = 0$, uncurved bar complex).  The semiclassical limit $c \to \infty$ ($\epsilon_2 \to 0$ in AGT, $b^2 = \epsilon_1/\epsilon_2 \to \infty$) gives Poisson brackets on the Hitchin base ($\kappa \to \infty$, strongly curved $A_\infty$ structure).  These are distinct limits: $c = 0$ is uncurved, $c \to \infty$ is maximally curved.
\end{example}

%================================================================
% SECTION: W-ALGEBRAS IN BOTH FRAMEWORKS
%================================================================

\section{$\mathcal{W}$-algebras in both frameworks}
% label removed: sec:w-algebras-unifying

\subsection{$\mathcal{W}$-algebras from 2D CFT perspective}

\begin{definition}[\texorpdfstring{$\mathcal{W}$}{W}-algebra (CFT definition)]% label removed: def:w-algebra-cft
Following Zamolodchikov \cite{Zamolodchikov} and Fateev--Lukyanov \cite{FL88}, a 
\emph{W-algebra} $\mathcal{W}$ is a vertex operator algebra containing:
\begin{enumerate}
\item Virasoro element $L$ (conformal weight 2)
\item Additional generators $W^{(s)}$ of conformal weights $s > 2$
\item Relations ensuring associativity of OPE
\end{enumerate}
\end{definition}

Standard examples: $\mathcal{W}_3$ has generators $L, W$ with $\text{wt}(L) = 2$, $\text{wt}(W) = 3$; $\mathcal{W}_N$ has generators of weights $2, 3, \ldots, N$; and $\mathcal{W}_{1+\infty}$ has infinitely many generators.

\subsection{$\mathcal{W}$-algebras from gauge theory perspective}

\begin{theorem}[Arakawa--Creutzig--Linshaw 2019; \ClaimStatusProvedElsewhere]% label removed: thm:w-from-gauge
\cite{ACL19} Let $G$ be a simple Lie group and $\rho: G \to GL(V)$ a 
representation. Consider the symplectic quotient (Higgs branch):
\[\mathcal{M}_H = \mu^{-1}(0) / G\]
where $\mu: T^*V \to \mathfrak{g}^*$ is the moment map.

Then:
\begin{enumerate}
\item The Higgs branch $\mathcal{M}_H$ carries a holomorphic symplectic structure

\item Quantization of functions $\mathcal{O}(\mathcal{M}_H)$ produces a vertex 
      algebra $V_H$
      
\item $V_H$ contains a W-algebra $\mathcal{W}_k(\mathfrak{g})$ at level $k$ 
      determined by the gauge coupling
      
\item This matches the W-algebra from coset construction: 
      \[\mathcal{W}_k(\mathfrak{g}) = \text{Com}(\mathfrak{g}_k, V_\rho)\]
\end{enumerate}
\end{theorem}


\subsection{Bar-cobar duality for $\mathcal{W}$-algebras}

\begin{conjecture}[\texorpdfstring{$\mathcal{W}$}{W}-algebra bar-cobar duality; \ClaimStatusConjectured]% label removed: conj:w-algebra-bar-cobar
Let $\mathcal{W}_k(\mathfrak{g})$ be a W-algebra (from either construction). Then:
\begin{enumerate}
\item \textup{[\ClaimStatusProvedHere]}
The chiral envelope $\mathcal{A}_{\mathcal{W}}$ admits geometric bar
      construction:
      \[\bar{B}^{\text{geom}}(\mathcal{A}_{\mathcal{W}}) = \bigoplus_{n \geq 0}
      \Gamma\left(\overline{C}_{n+1}(X), \mathcal{A}_{\mathcal{W}}^{\boxtimes(n+1)}
      \otimes \Omega^n_{\log}\right)\]

\item At generic level $k$ \textup{[\ClaimStatusProvedHere]};
      at admissible levels \textup{[\ClaimStatusConjectured]}:
      $\mathcal{W}_k(\mathfrak{g})$ is \emph{Koszul} and has a chiral
      Koszul dual coalgebra $\mathcal{A}_{\mathcal{W}}^!$

\item Conditional on~\textup{(2)}: the bar and cobar complexes are quasi-inverse:
      \[\Omega(\bar{B}(\mathcal{A}_{\mathcal{W}})) \simeq \mathcal{A}_{\mathcal{W}}\]
      At generic $k$ this is \textup{[\ClaimStatusProvedHere]};
      at admissible levels \textup{[\ClaimStatusConjectured]}

\item \textup{[\ClaimStatusProvedHere]}
All structures (Virasoro, W-currents, OPE) have geometric realization via
      configuration spaces
\end{enumerate}
\end{conjecture}

\begin{proof}[Proof of items \textup{(1)} and \textup{(4)}]
Item~(1): the geometric bar construction applies to any augmented chiral algebra
(Definition~\ref*{V1-def:geometric-bar}); the bar differential satisfies
$d^2 = 0$ by Theorem~\ref*{V1-thm:genus-induction-strict}.  For W-algebras,
the chiral envelope exists by the Beilinson--Drinfeld functor
(Theorem~3.7.11 of \cite{BD04}).

Item~(4): OPE data is encoded by residues at collision divisors of the
FM compactification (Definition~\ref*{V1-def:bar-differential-complete}).
For W-algebra currents $W^{(s)}(z)$, the OPE poles of orders
$\leq 2s$ correspond to the stratification of $\overline{C}_2(X)$
along the diagonal, with each pole order contributing to a specific
component of $d_{\mathrm{bar}}$.

Item~(3) follows from
Theorem~\ref*{V1-thm:bar-cobar-isomorphism-main} once the Koszul property~(2)
is established.
\end{proof}

\begin{proof}[Proof of items \textup{(2)} and \textup{(3)} at generic~$k$]
Let $k$ be generic (away from admissible and critical levels).

\emph{Step~1: Koszul property via weight spectral sequence.}
At generic~$k$, the W-algebra $\mathcal{W}_k(\mathfrak{g})$ is simple
as a vertex algebra (Arakawa \cite{Ara07}, Theorem~4.1).  Simplicity
implies that Verma modules are irreducible
(Proposition~\ref*{V1-prop:generic-irreducibility}; the Shapovalov form
is non-degenerate at each weight), so
$\mathrm{Ext}^q_{\mathcal{O}_k}(V(h_1), V(h_2)) = 0$ for $q > 0$
and all pairs of distinct Verma modules.

The weight filtration on the bar complex induces a spectral sequence
$E_2^{p,q} \Rightarrow \ChirHoch^{p+q}(\mathcal{W}_k(\mathfrak{g}))$.
The Ext-vanishing forces $E_2^{p,q} = 0$ for $q > 0$, so the spectral
sequence collapses at~$E_2$.  This $E_2$-degeneration is equivalent to
the Koszul property: the bar resolution is linear
(Adams-graded components are acyclic in positive internal degree).

\emph{Step~2: Identification of the Koszul dual.}
By Theorem~\ref*{V1-thm:w-algebra-hochschild} (\ClaimStatusProvedHere),
the bar complex cohomology at generic~$k$ is:
\[
\ChirHoch^*(\mathcal{W}_k(\mathfrak{g}))
\;=\; \mathbb{C}[\Theta_1, \ldots, \Theta_r],
\qquad \deg \Theta_i = h_i = m_i + 1
\]
where $m_1, \ldots, m_r$ are the exponents of~$\mathfrak{g}$.
The Koszul dual coalgebra $\mathcal{A}_{\mathcal{W}}^!$ is determined by
this cohomology: it is the cofree coalgebra cogenerated by the classes
$\Theta_i^*$ dual to the $\Theta_i$.

\emph{Step~3: Bar-cobar quasi-isomorphism.}
With the Koszul property established,
Theorem~\ref*{V1-thm:bar-cobar-isomorphism-main} gives the
quasi-isomorphism
$\Omega(\bar{B}(\mathcal{A}_{\mathcal{W}})) \simeq
\mathcal{A}_{\mathcal{W}}$.
\end{proof}


\begin{example}[Explicit: \texorpdfstring{$\mathcal{W}_3$}{W_3} at \texorpdfstring{$c = -2$}{c = -2}]% label removed: ex:w3-c-minus-2-explicit
For the $\mathcal{W}_3$ algebra at central charge $c = -2$:

\emph{Generators.}
\begin{align*}
L &= \text{energy-momentum tensor, conformal weight 2} \\
W &= \text{W-current, conformal weight 3}
\end{align*}

\emph{OPE.}
\begin{align*}
L(z)L(w) &\sim \frac{c/2}{(z-w)^4} + \frac{2L(w)}{(z-w)^2} + \frac{\partial L(w)}{z-w} \quad (c = -2) \\
L(z)W(w) &\sim \frac{3W(w)}{(z-w)^2} + \frac{\partial W(w)}{z-w} \\
W(z)W(w) &\sim \frac{c_{33}}{(z-w)^6} + \frac{2L(w)}{(z-w)^4} + \cdots
\end{align*}
where $c_{33} = c/3 = -2/3$ by the standard Zamolodchikov normalization.

\emph{Chiral algebra presentation.}
\[\mathcal{A}_{\mathcal{W}_3} = \text{Free}_{\mathcal{D}}(\mathbb{C}L \oplus \mathbb{C}W) 
/ \text{(W-algebra relations)}\]

\emph{Bar complex degree~2.}
\[\bar{B}^2(\mathcal{A}_{\mathcal{W}_3}) = \Gamma\left(\overline{C}_3(X), 
\mathcal{A}_{\mathcal{W}_3}^{\boxtimes 3} \otimes \Omega^\bullet\right)\]

Elements are represented by integrals:
\[\int_{\overline{C}_3(X)} f(z_1, z_2, z_3) \wedge d\log(z_1-z_2) \wedge d\log(z_2-z_3)\]
where $f$ is a section of $\mathcal{A}_{\mathcal{W}_3}^{\boxtimes 3}$.

\emph{Differential action.}
\begin{align*}
d_{\text{mult}}(f) &= \text{Res}_{z_1 = z_2}[f] + \text{Res}_{z_2 = z_3}[f] \\
&\quad + \text{Res}_{z_1 = z_3}[f] \quad \text{(collisions)} \\
d_{\text{internal}}(f) &= d_{\mathcal{W}}(f) \quad \text{(internal differential)} \\
d_{\text{extend}}(f) &= \text{extension across boundary divisors}
\end{align*}

\emph{Koszul dual coalgebra.} At $c = -2$, we conjecture (building on Arakawa's rationality result~\cite{Ara07}) that $\mathcal{W}_3$ is Koszul, with coalgebra dual $\mathcal{C}_{\mathcal{W}_3}$ given by:
\[\mathcal{C}_{\mathcal{W}_3} \subset \text{Cofree}_{\mathcal{D}}(s^{-1}\mathbb{C}L^* \oplus s^{-1}\mathbb{C}W^*)\]
the sub-coalgebra cogenerated by the orthogonal complement of the relations, where $s^{-1}$ denotes desuspension.
\end{example}

%================================================================
% SECTION: MATHEMATICAL BRIDGES
%================================================================

\section{AGT correspondence and bar-cobar duality}
% label removed: sec:mathematical-bridges



\subsection{AGT correspondence via bar-cobar}

\begin{theorem}[AGT 2D side: bar complex = semi-infinite complex; \ClaimStatusProvedHere]
% label removed: thm:agt-2d-bar
\index{AGT correspondence!2D side}
Let $\mathcal{W}_k(\fg)$ be a W-algebra at generic level~$k$.  Then $H^\bullet(\barB^{\mathrm{ch}}(\mathcal{W}_k(\fg))) \cong H^{\frac{\infty}{2}+\bullet}(\mathcal{W}_k(\fg))$ (Theorem~\ref*{V1-thm:bar-semi-infinite-w}), the anomaly duality $\kappa(\mathcal{W}_k(\fg)) + \kappa(\mathcal{W}_{k'}(\fg^\vee)) = 0$ holds for Feigin--Frenkel dual pairs (Corollary~\ref*{V1-cor:anomaly-duality-w}), and chiral homology is computed by configuration space integrals on FM compactifications (Theorem~\ref*{V1-thm:bar-NAP-homology}).
\end{theorem}

\begin{conjecture}[AGT 4D--2D bridge via bar-cobar; \ClaimStatusConjectured]% label removed: conj:agt-bar-cobar
\index{AGT correspondence}
The Alday--Gaiotto--Tachikawa correspondence \cite{AGT09} admits a
bar-cobar formulation: the Costello--Li holomorphic twist of 4D
$\mathcal{N}=2$ gauge theory with gauge group~$G$ produces a
factorization algebra $\mathcal{F}_G$ on a curve~$\Sigma$, and
this factorization algebra is equivalent to $\mathcal{W}_k(\fg)$:

\begin{center}
\begin{tikzcd}
\text{4D } \mathcal{N}=2 \text{ gauge partition function} \ar[d, "\text{CL twist}"] \ar[r, "\text{AGT}"]
& \text{2D Liouville/Toda CFT correlation} \ar[d, "\text{chiral envelope}"] \\
\text{2D factorization algebra } \mathcal{F}_G \ar[r, "\simeq"]
& \text{W-algebra } \mathcal{W}_k(\mathfrak{g})
\end{tikzcd}
\end{center}

Moreover:
\begin{enumerate}
\item The 4D instanton partition function = W-algebra conformal blocks
\item The 4D Coulomb branch parameters = W-algebra momenta
\item The BV complex of the 4D theory computes the same homology as
  the bar complex of the 2D chiral algebra:
      \[H_\bullet^{\text{ch}}(X, \mathcal{W}_k(\mathfrak{g})) = H_\bullet^{\text{BV}}(\mathcal{F}_G)\]
\end{enumerate}
\end{conjecture}

\begin{remark}[Scope]
The 2D side is proved (Theorem~\ref{thm:agt-2d-bar}); the 4D--2D bridge is proved in specific cases \cite{SV13,MO19}.  MC4 and MC5 are both proved in Volume~I (Chapter~\ref*{V1-chap:concordance}); the remaining obstruction is the 4D--2D bridge for general gauge groups.
\end{remark}


%================================================================
% SECTION: OPEN QUESTIONS
%================================================================

\section{Open questions and future directions}
% label removed: sec:open-questions-holo-topo

\begin{theorem}[\texorpdfstring{$E_n$}{E_n} Koszul duality on \texorpdfstring{$n$}{n}-manifolds;
\ClaimStatusProvedElsewhere]% label removed: thm:en-koszul
\index{En algebra@$E_n$ algebra!Koszul duality}
\index{Koszul duality!higher-dimensional}
For an $E_n$-algebra $\mathcal{A}$ on a smooth oriented $n$-manifold $M$,
there is a bar-cobar adjunction
\begin{equation}% label removed: eq:en-bar-cobar
\barB_{E_n}: E_n\text{-}\mathrm{Alg}(M)
\rightleftarrows
E_n\text{-}\mathrm{CoAlg}(M) :\Omega_{E_n}
\end{equation}
realized by configuration space integrals on the Axelrod--Singer
compactification $\overline{C}_k(M)$ of $k$-point configurations in $M$.
The propagator is the Green kernel of the Hodge Laplacian on $M$, and
the Arnold relations are replaced by the Totaro relations
\cite{Totaro96} for $H^*(\overline{C}_k(\mathbb{R}^n))$.

For $n = 1$, this recovers the classical (non-chiral) bar-cobar
adjunction for $A_\infty$-algebras.  For $n = 2$ on a Riemann surface,
it recovers the chiral bar-cobar adjunction of
Theorem~\textup{\ref*{V1-thm:bar-cobar-isomorphism-main}}.
The Poincar\'e--Koszul duality of Ayala--Francis~\cite{AF15}
\textup{(}Theorem~7.8\textup{)} provides the abstract framework;
our propagator formulas should supply explicit chain-level
representatives.

\emph{Scope.}  For $n \geq 3$, the configuration space integrals involve
the Axelrod--Singer propagator (which is not holomorphic), and the
combinatorics of boundary strata are governed by graph complexes
(Kontsevich \cite{Kon94}).  The key missing ingredient is an
explicit form-level presentation of the $E_n$ bar differential in terms
of iterated integrals over $\overline{C}_k(M)$.  The $n = 3$ case
with $M = S^3$ connects to the Chern--Simons perturbative expansion and
Vassiliev invariants (see Remark~\textup{\ref{rem:vassiliev}} below).

\emph{Resolution.}  All four components are now
established: see Theorem~\textup{\ref*{V1-thm:en-koszul-duality}}, which
assembles the proofs of Totaro, Lambrechts--Voli\'c, Knudsen, and
our $n = 2$ recovery.
\end{theorem}


\begin{remark}[Vassiliev invariants from Feynman transform]
% label removed: rem:vassiliev
\index{Vassiliev invariants!from bar complex}
The Prism Principle $\barB^{\mathrm{full}} \simeq \mathrm{FT}_{\mathrm{Mod}}$ (Theorem~\ref*{V1-thm:prism-higher-genus}(ii)) and Getzler--Kapranov \cite{GeK98} suggest (\ClaimStatusConjectured) that the bar complex of the Chern--Simons factorization algebra on $S^3$ produces universal Vassiliev invariants, recovering the Kontsevich integral via specialization to the real locus.  The main gap is the passage $\overline{C}_n(X) \to \overline{C}_n(S^1)$.
\end{remark}

